%% encoding: UTF-8
%% 首先读取模板的配置
%%%%%%%%%%%%%%%%%%%%%%%%%%%%%%%设置版面控制%%%%%%%%%%%%%%%%%%%%%%%%%%%%%%%%%%
\documentclass[a4paper, 12pt, notitlepage, oneside , twoside ]{article}
\usepackage[colorlinks, linkcolor=red, anchorcolor=blue, citecolor=green, bookmarks]{hyperref}
\hypersetup{xetex}
% 设置版面，纸张空距
\usepackage[papersize={210mm,285mm},headheight=2cm,top=2.05cm,headsep=0cm,bottom=0.82cm,footskip=.2cm,
left=0cm,right=2.06cm
]{geometry}
\usepackage{marvosym}
\usepackage{makecell,multirow}
\usepackage{wrapfig}
\usepackage{setspace}

% 添加首行缩进，两个字符
\usepackage{indentfirst}
\setlength{\parindent}{2em}

\usepackage{amssymb}
\usepackage{amsmath}
\usepackage{fancyhdr}
\pagestyle{fancy}
%%%%%%%%%%%%%%%%%%%%%%%%%%%%%设置字体%%%%%%%%%%%%%%%%%%%%%%%%%%%%%
\usepackage[no-math]{fontspec} %提供字体选择命令
\usepackage{xunicode,xltxtra}
\usepackage[CJKchecksingle, CJKnumber]{xeCJK}         % 用于字体
\usepackage{CJK}
\setCJKfamilyfont{kai}[BoldFont=Adobe Heiti Std]{Adobe Kaiti Std}
\setCJKfamilyfont{song}[BoldFont=Adobe Heiti Std]{Adobe Song Std}
\setCJKfamilyfont{hei}[BoldFont=Adobe Heiti Std]{Adobe Heiti Std}
\setCJKmainfont[BoldFont={FZLanTingHei-DB-GBK}]{FZLanTingHei-L-GBK}
\setmainfont{FZLanTingHei-L-GBK}  % 设置中文/粗体/斜体/主要字体
\punctstyle{hangmobanjiao}        % 设置行末半角
\newfontfamily\times{Times New Roman}
\newCJKfontfamily\kaishu{Adobe Kaiti Std}
\setCJKfamilyfont{zhyh}{Microsoft YaHei}
\newcommand*{\heiti}{\CJKfamily{zhyh}}
\setCJKfamilyfont{zhltxh}{FZLanTingHei-L-GBK}
\newcommand*{\lantxh}{\CJKfamily{zhltxh}}
\newfontfamily\arial{Arial}
\newcommand{\barial}{\bf\arial}
%\setmainfont[Mapping=tex-text]{Times New Roman}         % 设置衬线英文字体
\setsansfont[Mapping=tex-text]{FZLanTingHei-L-GBK}         % 设置无衬线英文字体
%\setmonofont{Courier New}                % 设置等款英文字体
\newcommand{\chuhao}{\fontsize{42pt}{\baselineskip}\selectfont}
\newcommand{\xiaochuhao}{\fontsize{36pt}{\baselineskip}\selectfont}
\newcommand{\yihao}{\fontsize{28pt}{\baselineskip}\selectfont}
\newcommand{\erhao}{\fontsize{21pt}{\baselineskip}\selectfont}
\newcommand{\xiaoerhao}{\fontsize{19pt}{\baselineskip}\selectfont}
\newcommand{\sanhao}{\fontsize{15.75pt}{\baselineskip}\selectfont}
\newcommand{\sihao}{\fontsize{14.01pt}{\baselineskip}\selectfont}
\newcommand{\xiaosihao}{\fontsize{11pt}{\baselineskip}\selectfont}
\newcommand{\wuhao}{\fontsize{10.04pt}{\baselineskip}\selectfont}
\newcommand{\xiaowuhao}{\fontsize{9.03pt}{15pt}\selectfont}
\newcommand{\liuhao}{\fontsize{8.53pt}{\baselineskip}\selectfont}
\newcommand{\qihao}{\fontsize{5.25pt}{\baselineskip}\selectfont}
\newcommand{\RNum}[1]{\uppercase\expandafter{\romannumeral #1\relax}}
\newcommand{\rnum}[1]{\lowercase\expandafter{\romannumeral #1\relax}}

\usepackage{textgreek}
\usepackage{pdflscape}
%%%%%%%%%%%%%%%%%%%%%%%%%%%%%设置图像%%%%%%%%%%%%%%%%%%%%%%%%%%%%%
\usepackage[usenames,dvipsnames,svgnames,table,cmyk]{xcolor}
\usepackage{graphicx}          % 用于图像编号
\graphicspath{{_drugRisk/}{_logo/}{_phenoDegree/}{disease_pic/}{drug_pic/}{pheno_pic/}{_barcode/}{PDF/}{cores/}}
\usepackage{tikz}
\usepackage{verbatim}
\usepackage{animate}
\usepackage{subfig}
\usepackage{wrapfig}
\usepackage{pgfplots}
\usepackage{wallpaper}%%%%%%%%%%%%%%%%背景%%%%%%%%%%%%%%%%%%
\usepackage{color}
\definecolor{gray2}{cmyk}{0.00,0.00,0.00,0.70}
\definecolor{blackdesc}{cmyk}{0, 0, 0, 1}	%黑色字体
\definecolor{titleblue}{cmyk}{1, 0.3, 0, 0}	%标题蓝色，线条
\definecolor{titledarkgreen}{cmyk}{0.8, 0, 0.65, 0}	%标题深绿，线条
\definecolor{titlegrassgreen}{cmyk}{0.5, 0, 1, 0}	%标题草绿，线条
\definecolor{tablegreen}{cmyk}{0.9, 0, 0.9, 0}  %table head color

%%%%%%%%%%%%%%%%%%%%%%%%%%mdframed%%%%%%%%%%%%%%%%%%%%
\usepackage[framemethod=tikz]{mdframed}
\usetikzlibrary{patterns}
\mdfdefinestyle{titlesty}{%
linecolor=platinum,
middlelinewidth=1.5pt,
frametitlerule=false,%
apptotikzsetting={\tikzset{
    mdfbackground/.append style ={
    top color=white,
    bottom color=white
    },
    mdfframetitlebackground/.append style={
    shade,
    top color=white ,
    bottom color=white
    }
}
},
frametitlerulewidth=0pt,
innertopmargin=\topskip,
topline=false,
leftline=false,
rightline=false,
bottomline=false,
frametitlealignment=\raggedright,
}
\usepackage{fancyhdr}
\pagestyle{fancy}
\fancyhf{}
\renewcommand{\headrulewidth}{0pt}
\renewcommand{\footrulewidth}{0pt}

%%%%%%%%%%%%%%%产品推荐框格式%%%%%%%%%%%%%%%%%%%%%%%%%%%%
\usepackage{paralist}
\usepackage{float,longtable,hhline}
\usepackage{ulem}

\definecolor{topcolor}{cmyk}{0.28,0,1,0}
\definecolor{edgecolor}{cmyk}{0.5,0,1,0}
\definecolor{markcolor}{cmyk}{0,1,1,0.2}

\definecolor{black100}{cmyk}{0.00,0.00,0.00,1.00}
\definecolor{black70}{cmyk}{0.00,0.00,0.00,0.70}

\usepackage[many]{tcolorbox}
\tcbuselibrary{skins}

\newtcolorbox{mybox}[1]{hbox boxed title,
  enhanced,frame hidden,colback=white,title=Sharp corners,
  attach boxed title to top left={yshift=-3mm,xshift=5mm},
  boxed title style={colframe=edgecolor,colback=edgecolor,size=small},
  borderline={0.5mm}{0mm}{edgecolor,dotted},
  title={#1}
}

%%%%%%%%%%%%%%%页码修改[RO]居右[CO]居中[LO]居左%%%%%%%%%%%
\begin{document}

\hspace*{12.5mm}\includegraphics[scale=1]{logo3.pdf}

\vspace*{11mm}

\begin{center}
{\bf\xiaoerhao\color{black70} \qquad\qquad 未来优态肠道调理建议}
\end{center}

\hspace*{12mm}
\begin{minipage}[b]{.89\hsize}
\vspace*{3mm}

\begin{spacing}{1.5}
\color{black70}{\fontsize{13pt}{13pt}\selectfont 检测编号：GI00001 \qquad 报告日期：2020-03-31 \\}
\color{black70}{\fontsize{13pt}{17pt}\selectfont 尊敬的\uline{\qquad 荔湾 \qquad} 
男士：
}
\end{spacing}

\begin{spacing}{1.5}
\setlength{\parindent}{2em}\color{black70}\fontsize{9.5pt}{12pt}\selectfont {
{\fontsize{13pt}{13pt}\selectfont 您好！}

感谢您选择常道康{\textsuperscript{\qihao{TM}}}肠道微生态检测套餐服务，我们将为您提供最科学的肠道微生态调理方案。

肠道微生态与我们的健康息息相关。工作压力、饮食不规律、食品安全、雾霾、滥用抗生素等都会导致肠道微生态失调，而80{\%}的健康问题与肠道微生态紊乱有关。

微生态制剂是调理肠道微生态达到肠道健康最直接有效的方式。微生态制剂包括益生元、益生菌等，可以促进益生菌的生长，抑制有害菌的生长，改善肠道微生态平衡，可以双向调节便秘腹泻、改善肥胖、增强机体免疫力等。

未来优态肠道调理产品根据肠道微生态基因检测结果，量身定制调理产品、剂量、周期等进行精准调理。调理周期结束后建议进行常道康{\textsuperscript{\qihao{TM}}}复查，准确评估干预效果及干预方案的调整，形成精准健康管理闭环。
}

\vspace*{3mm}
\end{spacing}

\begin{spacing}{1.5}
\begin{mybox}{根据您本次肠道微生态检测结果，建议您采用以下肠道调理方案}
\color{black70}\fontsize{9.5pt}{12pt}\selectfont
\begin{tabular}{m{5cm}<{\centering}m{8cm}<{\centering}}
  \parbox[c][0.6cm]{\hsize}{\color{black70}\bf 肠道调理产品：} &
  \parbox[c][0.6cm]{\hsize}{\color{black70}\bf 推荐用量：} \\
            \parbox[c][0.6cm]{\hsize}{\color{black70} 产品 ：\uline{\qquad 未来优态UA \qquad}，} &
      \parbox[c][0.6cm]{\hsize}{\color{black70} 用法：1天\uline{\quad 2 \quad}次，1次\uline{\quad 1 \quad}袋，餐\uline{\quad 后 \quad}服用} \\
      \end{tabular}

\vspace*{2mm}
  {~~ 建议调理}
            {1 - 3 个月。}
      {（第一个月干预产品预计使用量：}
      {2 盒 UA }
  {）}\\
\fontsize{8.5pt}{12pt}\selectfont {~~{\raisebox{1pt}{\color{edgecolor}$\bullet$}}\! 本方案为推荐使用方案，具体服用量和调理周期可根据自身实际情况和服用效果进行调整。若有问题请咨询健康顾问。}
\end{mybox}
\end{spacing}
\vspace*{-3mm}
\begin{spacing}{1.5}
            \noindent\color{black70}\fontsize{8.5pt}{12pt}\selectfont {{\raisebox{1pt}{\color{markcolor}$\divideontimes$}}\! 温馨提示 ：益生菌固体饮料为活菌制剂，建议用37℃以下温水冲调并及时服用，以免有益菌失活而影响干预效果}\\
      \end{spacing}

\vspace*{2mm}
\hspace*{12cm}\includegraphics[scale=0.6]{erweima.pdf}

\vspace*{2mm}
\hfill\color{black70}\fontsize{13pt}{13pt}\selectfont 了解肠道健康资讯，请关注“常道康”
\end{minipage}

\fancyfoot[L]{\includegraphics{footer.pdf}}

\end{document}
